% !PRINTED 670
\weberpage{670}
%<<CONT
worin $a, b, c, d$ Konstanten sind, deren Determinante $ad - bc$ von Null verschieden ist, so sind die Verzweigungspolygone $\mathfrak{Z}_z$ und $\mathfrak{Z}_{z'}$ identisch, und daher ist auch $w_z = w_{z'}$, denn wenn in einem Punkte $z-z_0$ oder $1:z$ Null in der $e$ten Ordnung wird, so ist in demselben Punkte auch

\[
z' - z_0' = \dfrac{(ad-bc)(z-z_0)}{(cz+d)(cz_0+d)}
\]

oder wenn $z_0$ unendlich, also $z_0' = a:c$ ist:

\[
z' - z_0' = \dfrac{-(ad-bc)}{c(cz+d)},
\]

und wenn $z_0' = \infty$, also $cz_0+d = 0$ ist:

\[
\dfrac{1}{z'} = \dfrac{cz+d}{az+b}
\]

ebenfalls unendlich klein in der $c$ten Ordnung.
% [sic] printed $c$ten Ordnung, should be $e$ten: the order $e$ in which $z-z_0$ or $1:z$ becomes zero (stated above); $c$ is one of the constants $a, b, c, d$

Wenden wir dies auf $z' = 1:z$ an, so erhalten wir die Verzweigungspunkte, in denen $z$ einen endlichen Wert hat, aus der Körperdiskriminante $D_z$ nach § 184, (24), (28). Es fehlen dann noch die Verzweigungspunkte, in denen $z=\infty$, also $z'=0$ wird, und diese ergeben sich in gleicher Weise aus den verschwindenden Wurzeln der Körperdiskriminante $D_{z'}$. Demnach haben wir den Satz:

5. \gesperrt{Die Verzweigungszahl $w_z$ ist gleich dem Grade der Körperdiskriminante $D_z$, vermehrt um die Anzahl der verschwindenden Wurzeln von $D_{z'}$. ($z' = 1:z$).}

\begin{center}
\textbf{§ 191. Polygonquotienten und Polygonklassen.}
\end{center}

Eine Funktion $\eta$ in $\Omega$ hat nur in einer endlichen Zahl von Punkten eine von Null verschiedene Ordnungszahl. Die Summe der positiven Ordnungszahlen ist ebenso groß wie die Summe der negativen, nämlich gleich der Ordnung der Funktion $\eta$ (§ 189, 4.). Sind die Ordnungszahlen von $\eta$ für jeden Punkt $\mathfrak{P}$ bekannt, so ist dadurch die Funktion $\eta$ bis auf einen konstanten Faktor bestimmt. Denn wenn eine zweite Funktion $\eta'$ überall dieselben Ordnungszahlen hat, so hat $\eta:\eta'$ überall die Ordnungszahl $0$ und ist daher konstant (§ 189, 4.).

Wir bilden ein Polygon $\mathfrak{A}$, in das wir jeden Punkt, in dem $\eta$ eine positive Ordnungszahl hat, so oft aufnehmen, als diese Ordnungszahl angibt, und ein zweites Polygon $\mathfrak{B}$, in das wir in
%CONT>>
